\section{Introduction}

A genuinely autonomous knowledge system cannot be a mere pipeline of
inference steps; it must survive for months or years, gracefully handle
resource fluctuations, and justify its own computational expenditures in
terms of information gain.  These requirements demand an infrastructure
that is \emph{self‑monitoring} and \emph{self‑regulating}---qualities
that are conspicuously absent in most AI frameworks, which treat hardware
as an infinite resource and failures as rare exceptions.

The \Apeiron{} Framework~\cite{beniers2026apeiron} aims to build an AI
that constructs knowledge from first principles without human‑provided
categories.  Its seventeen‑layer design moves from irreducible observables
to pluri‑versal world‑making, but every layer rests on a shared
foundational substrate that must guarantee responsiveness, auditability, and
ethical resource consumption.  This paper describes that substrate: the
\Core{}.

The \Core{} consists of four interoperating modules:
\begin{enumerate}
    \item \textbf{EventBus} -- a publish‑subscribe message bus with
          persistent storage, dead‑letter queues, and replay capability,
          enabling decoupled communication across all layers.
    \item \textbf{AdaptiveThresholds} -- a self‑calibrating threshold
          manager that adjusts safety margins, stability criteria, and
          sensitivity parameters in response to the system's current
          complexity and the number of unresolved gaps.
    \item \textbf{ChaosDetector} -- a runtime monitor that classifies the
          system state (stable, oscillating, diverging, chaotic) and triggers
          programmed interventions, using Lyapunov exponent estimation,
          Kalman filtering, and a circuit breaker.
    \item \textbf{ThermodynamicCost} -- an energy‑aware cost model that
          tracks power consumption via hardware telemetry (CPU, GPU),
          optionally queries real‑time CO\(_2\) intensity, and approves
          or rejects computational structures based on their
          information‑value‑per‑Joule.
\end{enumerate}

Together, these modules provide the operational substrate that enables
the formal verification and self‑proving capabilities described in our
companion work on multi‑prover atomicity certification~\cite{beniers2026multi}.
Without a self‑regulating infrastructure, the guarantees obtained from
the theorem provers would be confined to isolated experiments; the \Core{}
embeds them into a continuously running, resilient agent.

The contribution of this paper is the \emph{integrated design} of these
four modules and an empirical demonstration that they can be composed into
a coherent self‑regulation loop.  We are not the first to propose
event‑driven architectures, adaptive thresholds, chaos engineering, or
energy‑aware computing.  Our contribution lies in combining them into a
single, lightweight infrastructure that is purpose‑built for autonomous
knowledge systems, is fully implemented, and is accompanied by reproducible
benchmarks.

The rest of the paper is organised as follows.  Section~2 surveys related
work.  Section~3 gives a high‑level architectural overview.  Sections~4--7
detail the event bus, adaptive thresholds, chaos detector, and
thermodynamic cost model, respectively.  Section~8 presents experimental
results, Section~9 discusses limitations and future directions, and
Section~10 concludes.